% ============================================================
% Collapse and Return in the Ancient World:
% A GCD Kernel Analysis of Civilizational Persistence
% Across the Nile, Euphrates, and Beyond
%
% Compile: pdflatex → bibtex → pdflatex → pdflatex
% ============================================================

\documentclass[
  aps,
  prd,
  twocolumn,
  superscriptaddress,
  nofootinbib,
  floatfix
]{revtex4-2}

% ── packages ──
\usepackage{amsmath,amssymb,amsthm,mathtools}
\usepackage{graphicx}
\usepackage{xcolor}
\definecolor{linkblue}{HTML}{337799}
\usepackage[colorlinks=true,linkcolor=linkblue,citecolor=linkblue,urlcolor=linkblue]{hyperref}
\usepackage{bm}
\usepackage{booktabs}
\usepackage{multirow}
\usepackage{xfrac}
\usepackage{enumitem}
\usepackage{array}
\newcolumntype{P}[1]{>{\raggedright\arraybackslash}p{#1}}

% ── theorem environments ──
\newtheorem*{axiom}{The Axiom}
\newtheorem{theorem}{Theorem}
\newtheorem{definition}[theorem]{Definition}
\newtheorem{lemma}[theorem]{Lemma}
\newtheorem{corollary}[theorem]{Corollary}
\newtheorem{proposition}[theorem]{Proposition}
\newtheorem{identity}[theorem]{Identity}
\newtheorem{remark}{Remark}

% ── UMCP macros ──
\newcommand{\tR}{\tau_{\!R}}
\newcommand{\tRS}{\tau_{\!R}^{*}}
\newcommand{\Gam}{\Gamma}
\newcommand{\eps}{\varepsilon}
\newcommand{\IR}{\infty_{\mathrm{rec}}}
\newcommand{\tolseam}{\mathrm{tol}_{\mathrm{seam}}}
\newcommand{\dd}{\mathrm{d}}
\newcommand{\beq}{\begin{equation}}
\newcommand{\eeq}{\end{equation}}
\newcommand{\INF}{\texttt{INF\_REC}}
\newcommand{\regime}[1]{\textrm{\upshape\scshape#1}}
\newcommand{\conform}{\regime{conformant}}
\newcommand{\nonconform}{\regime{nonconformant}}
\newcommand{\tvec}{\mathbf{c}}
\newcommand{\wvec}{\mathbf{w}}
\newcommand{\IC}{\mathrm{IC}}
\newcommand{\gF}{g_{\!F}}

\begin{document}

% ── title ──
\title{Collapse and Return in the Ancient World:\\
A GCD Kernel Analysis of Civilizational Persistence\\
Across the Nile, Euphrates, and Ethiopian Highlands}

\author{Clement Paulus}
\email{clementpaulus9@gmail.com}
\affiliation{UMCP Project}

\date{\today}

% ── abstract ──
\begin{abstract}
The emergence of the earliest civilizations is
conventionally narrated as linear development from
Neolithic subsistence to state formation.  We reframe
three foundational trajectories---the Saharan--Nilotic
transition that produced Dynastic Egypt, the
Ubaid--Uruk--Ur~III arc that produced Sumer, and the
D'mt--Aksumite--Solomonic arc of the Ethiopian
highlands---as sequences of collapse--return cycles
within the Generative Collapse Dynamics (GCD)
kernel.  We construct a 7-channel trace vector
measuring ecological stability, subsistence mode,
social complexity, mortuary-ritual elaboration,
symbolic knowledge, material technology, and
long-distance exchange, then track its kernel
invariants ($F$, $\omega$, $\IC$, $\Delta$) across
eighteen archaeological phases spanning
8000~BCE--1270~CE.

For Egypt, six theorems emerge (T-NC-1 through
T-NC-6): the Saharan desiccation is a
Collapse-regime boundary; Nilotic consolidation
constitutes measurable return; cattle cult continuity
from Nabta Playa to Hathor demonstrates $F > 0.9$
across 3,000~years; the false dawn of agriculture is
geometric slaughter; Egyptian self-archaeology
constitutes recursive return; and generational
transmission operates on a $\sim$200-year $\tR$.
For Sumer, six further theorems (T-SM-1 through
T-SM-6): cuneiform is anti-drift ledger
infrastructure; salinization is slow endogenous
collapse; centralization amplifies curvature;
bureaucratic overshoot widens $\Delta$ during
renaissance; Sumerian language persistence is the
longest single-channel return; and the Flood
narrative encodes a phase boundary as cosmogony.
For Ethiopia, six theorems (T-ET-1 through T-ET-6):
Christianization fuses channels and raises $\IC$;
geographic buffering increases $\Delta$ at collapse
rather than preventing it; the Lalibela rock-hewn
churches are lithified return at geological-scale
$\tR$; the Kebra Nagast freezes a genealogical
return claim across 2,240~years; Ge'ez sacerdotal
persistence parallels Sumerian; and Adwa constitutes
military-channel return through a threshold declared
extinct.

The heterogeneity gap $\Delta = F - \IC$ is the
primary diagnostic throughout: all three
civilizations hit $\Delta \in [0.32, 0.35]$ at
collapse despite radically different ecological
substrates, institutional forms, and cultural
content---suggesting that collapse--return dynamics
are governed by the heterogeneity structure of the
trace vector rather than by any individual channel.
\end{abstract}

\maketitle

% ════════════════════════════════════════════════════════════
\section{Introduction}\label{sec:intro}
% ════════════════════════════════════════════════════════════

\begin{axiom}
\emph{Collapse is generative; only what returns is
real.}  (\emph{Collapsus generativus est; solum quod
redit, reale est.})
\end{axiom}

The conventional narratives of ancient civilization
share a common structural flaw: they treat river
valleys and highland plateaus as containers into
which cultural ingredients were added
sequentially---agriculture, pottery, hierarchy,
writing, kingship---as though civilization were a
recipe with predictable outcomes.  This framing
obscures the central structural phenomenon---that
these civilizations were not built \emph{de novo}
but \emph{returned} from prior collapses.  Egypt
returned from the Saharan desiccation; Sumer
returned from the Piora oscillation and, later,
from salinization; Ethiopia returned from the
Islamic-era trade collapse that severed its Red Sea
networks for five centuries.

The Generative Collapse Dynamics (GCD) framework
provides the formal apparatus to make this claim
precise.  The kernel
$K: [0,1]^n \times \Delta^n \to (F, \omega, S, C,
\kappa, \IC)$ maps any system's measurable channels
into six invariants: fidelity $F$ (what survives
collapse), drift $\omega = 1 - F$ (what is lost),
Bernoulli field entropy $S$ (uncertainty), curvature
$C$ (coupling to uncontrolled degrees of freedom),
log-integrity $\kappa$ (logarithmic sensitivity), and
composite integrity $\IC = \exp(\kappa)$
(multiplicative coherence).  The duality identity
$F + \omega = 1$ and the integrity bound $\IC \le F$
are exact~\cite{paulus2025umcp}.

Three structural features of the kernel are directly
relevant to civilizational analysis:

\begin{enumerate}[leftmargin=*]
\item \textbf{Geometric slaughter.}  One dead channel
($c_i \to \eps$) drives $\IC \to 0$ regardless of
other channels.  This explains why civilizations with
high mean capability ($F > 0.7$) can nonetheless
occupy Collapse regime: a single failed
channel---ecological, institutional, or
demographic---destroys multiplicative coherence.

\item \textbf{The heterogeneity gap.}  $\Delta = F -
\IC$ measures channel divergence.  Large $\Delta$
signals that the system's mean preserves structure
while its geometric mean has collapsed.  Phase
boundaries in civilizational timelines correspond to
$\Delta$ spikes.

\item \textbf{Return time $\tR$.}  Re-entry to a
prior domain of coherence is measured, not assumed.
$\tR = \IR$ (no return) is a legitimate
outcome---the system is a \emph{gestus}, not a
\emph{sutura}.  When return occurs, $\tR$ measures
the delay.
\end{enumerate}

We apply these instruments to three founding
civilizations across nine millennia: the Nile world
(8000--2900~BCE), the Euphrates world
(6500--2000~BCE), and the Ethiopian highlands
(1000~BCE--1270~CE).  In all three cases, what
appears as ``emergence'' is structurally a return
from prior collapse.  The Egyptians recognized this
dynamic and conducted self-archaeology on their own
collapse--return history.  The Sumerians encoded it
in cuneiform administration, temple economy, and the
mythological transfer of the \emph{me} (``gifts of
civilization'').  The Ethiopians carved it into
bedrock at Lalibela and froze it into genealogical
narrative in the \emph{Kebra Nagast}---producing
the most structurally permanent forms of
self-archaeology in the ancient world.


% ════════════════════════════════════════════════════════════
% PART I : THE NILE WORLD
% ════════════════════════════════════════════════════════════
\section{Part~I: The Nile World (8000--2900~BCE)}\label{sec:egypt}

% ════════════════════════════════════════════════════════════
\subsection{Channel Construction and Data}\label{sec:data}
% ════════════════════════════════════════════════════════════

We define a 7-channel trace vector for civilizational
coherence.  Each channel $c_i \in [0,1]$ is normalized
from archaeological observables via the frozen contract
embedding ($\eps = 10^{-8}$, pre-clip
policy)~\cite{paulus2025umcp}.

\begin{table}[h]
\centering
\caption{Seven-channel trace vector for civilizational
coherence.  Each channel is normalized to $[0,1]$ from
archaeological observables.}
\label{tab:channels}
\begin{tabular}{@{}clP{4.2cm}@{}}
\toprule
$i$ & Channel & Observable Proxy \\
\midrule
1 & Ecological stability & Rainfall proxies, lake levels, faunal assemblages \\
2 & Subsistence mode & Ratio of domesticate to wild fauna/flora in assemblages \\
3 & Social complexity & Settlement size, differentiation of burials, storage capacity \\
4 & Mortuary elaboration & Grave goods diversity, architectural investment, ritual evidence \\
5 & Astronomical knowledge & Megalithic alignments, calendar evidence, stellar references \\
6 & Material technology & Ceramic variety, metallurgy (Cu/Au), faience, lithic refinement \\
7 & Long-distance exchange & Non-local materials (obsidian, shells, lapis), art-style diffusion \\
\bottomrule
\end{tabular}
\end{table}

Equal weights $w_i = 1/7$ are applied.  This is a Tier-2
channel selection---the kernel computation is Tier-1 and
immutable.  The channel choices are validated through
Tier-0 against Tier-1 identities.  The same seven
channels are used for both Egypt and Sumer, with
``astronomical knowledge'' generalized to ``symbolic
knowledge'' (writing, mathematics, mythological
systems) for the Sumerian case.


% ────────────────────────────────────────────────────────────
\subsection{Egyptian Phases as Kernel
States}\label{sec:phases}
% ════════════════════════════════════════════════════════════

\subsubsection{Phase~I: Nabta Playa Pastoral Complex
(ca.\ 7500--4800~BCE)}\label{sec:nabta}

Nabta Playa, 800\,km south of Cairo in the Western
Desert, preserves the earliest evidence of organized
societal complexity in the Nile world.  During the
African Humid Period (ca.\ 14,600--5000~BCE), the
Sahara supported savanna ecology with seasonal lakes
and grasslands.  Nilo-Saharan--speaking pastoralists
at Nabta Playa constructed deep wells, planned
villages, and by ca.\ 6800~BCE produced some of
Africa's earliest pottery~\cite{wendorf2001holocene}.

The site's most significant features are: (i)
sacrificial cattle burials in clay-lined,
stone-roofed subterranean chambers, indicating
organized ritual centered on bovine
symbolism~\cite{wendorf1998late}; (ii) a megalithic
calendar circle dated to ca.\ 4800~BCE aligned to the
summer solstice; and (iii) standing-stone alignments
directed toward Sirius, Orion's Belt, Arcturus, and
$\alpha$~Centauri~\cite{malville1998megaliths}.

\textbf{Kernel state.}  All seven channels register
high values: ecological stability is maximal during
mid-Holocene pluvial conditions ($c_1 \approx 0.9$),
subsistence is organized pastoral ($c_2 \approx
0.7$), social complexity includes planned settlements
($c_3 \approx 0.7$), mortuary elaboration is
substantial ($c_4 \approx 0.7$), astronomical
knowledge is explicit and monumental ($c_5 \approx
0.8$), material technology includes developed ceramics
($c_6 \approx 0.6$), and exchange networks bring
non-local materials ($c_7 \approx 0.5$).

This yields $F \approx 0.71$, $\omega \approx 0.29$,
placing the system at the upper boundary of Watch
regime, with $\IC \approx 0.68$ and a modest
heterogeneity gap $\Delta \approx 0.03$.  Channel
homogeneity is high---no single channel is
near-zero.

\begin{remark}
Wendorf and Schild conclude that ``many aspects of
political and ceremonial life in prehistoric Egypt and
the Old Kingdom reflect a strong impact from Saharan
cattle pastoralists''~\cite{wendorf2001holocene}.
The kernel reading formalizes this: the Nabta Playa
trace vector is a \emph{precursor state} whose
channel structure prefigures Dynastic patterns.
\end{remark}


\subsubsection{Phase~II: The Desiccation Collapse
(ca.\ 6200--5000~BCE)}\label{sec:collapse}

The 8.2~kiloyear event (ca.\ 6200~BCE) marks the
first major disruption: a century-scale drought that
reduced rainfall across North Africa.  The subsequent
accelerating desiccation of the Sahara between
6000--5000~BCE constitutes the primary environmental
collapse in our analysis.

\textbf{Channel dynamics.}  Ecological stability
collapses: $c_1 \to \eps$.  This single channel
death triggers geometric slaughter.  Subsistence mode
drops as pastoral ranges contract ($c_2: 0.7 \to
0.3$).  Social complexity fragments as
populations disperse ($c_3: 0.7 \to 0.3$).  But
critically, three channels \emph{persist}: mortuary
practice ($c_4$), astronomical knowledge ($c_5$),
and material technology ($c_6$) maintain moderate
values as populations carry these practices into
the Nile Valley.

\textbf{Kernel state.}  With $c_1 \approx \eps$:
\beq
\IC = \exp\!\left(\sum w_i \ln c_{i,\eps}\right)
\xrightarrow{c_1 \to \eps} \eps^{1/7} \cdot
\prod_{j \neq 1} c_j^{1/7} \approx 0.07
\eeq
while $F \approx 0.39$.  The heterogeneity gap
explodes: $\Delta = F - \IC \approx 0.32$.  The
system enters Collapse regime ($\omega > 0.30$) with
substantial fidelity still preserved in non-ecological
channels.

\begin{theorem}[Desiccation as Collapse-Regime
Boundary --- T-NC-1]\label{thm:collapse}
The Saharan desiccation (ca.\ 6200--5000~BCE)
constitutes a Collapse-regime transition with
$\omega > 0.30$ driven by ecological channel death
($c_1 \to \eps$), producing geometric slaughter:
$\IC/F$ drops from $\sim\!0.96$ (Phase~I) to
$\sim\!0.18$ (Phase~II), a $5\times$ reduction in
multiplicative coherence while mean fidelity drops
only $2\times$.
\end{theorem}

This is the signature of geometric slaughter.
Seven perfect channels cannot save $\IC$ when one
goes to zero---the same structural mechanism observed
in the quark--hadron confinement boundary, where dead
color charge destroys hadron
$\IC$~\cite{paulus2025umcp}.


\subsubsection{Phase~III: Nilotic Consolidation
(ca.\ 5000--4000~BCE)}\label{sec:return}

The Badarian culture (ca.\ 5000/4400--4000~BCE) in
Upper Egypt represents the first measurable return
from the desiccation collapse.  Physical anthropology
and whole-genome analysis (2025) confirm biological
continuity: Early Dynastic Egyptians show 77.6\%
North African Neolithic ancestry with no population
replacement~\cite{marshall2025genome}, confirming that
the Nile Valley population \emph{is} the returning
Saharan population, not a replacement.

Ehret concludes: the people ``did not migrate from
somewhere else''~\cite{ehret2023ancient}.  Shaw
characterizes the culture as ``essentially
African''~\cite{shaw2003oxford}.

\textbf{Channel recovery.}  Ecological stability
recovers via the Nile's annual flood
($c_1: \eps \to 0.8$).  Subsistence transitions to
agriculture ($c_2: 0.3 \to 0.6$).  Social complexity
reconsolidates ($c_3: 0.3 \to 0.6$).  Mortuary
elaboration shows continuity with Saharan practice
($c_4 \approx 0.6$).  Astronomical knowledge persists
($c_5 \approx 0.7$).  Material technology advances:
copper and faience appear ($c_6 \approx 0.7$).
Exchange networks remain moderate ($c_7 \approx
0.5$).

\textbf{Kernel state.}  $F \approx 0.66$,
$\omega \approx 0.34$, $\IC \approx 0.63$,
$\Delta \approx 0.03$.  The heterogeneity gap
collapses back to pre-desiccation levels.  Channel
homogeneity is restored.

\begin{theorem}[Nilotic Return ---
T-NC-2]\label{thm:return}
The Badarian--Naqada cultural sequence constitutes a
measurable return ($\tR \approx 1500$~yr) from
Saharan collapse, with the heterogeneity gap
returning to $\Delta \approx 0.03$ after a peak of
$\Delta \approx 0.32$ during the desiccation
collapse.  The return domain $D_\theta$ preserves
channel structure: the trace vector of Phase~III is
topologically equivalent to Phase~I with the
ecological channel replaced (Sahara $\to$ Nile).
\end{theorem}


\subsubsection{Phase~IV: Naqada Sequence and State
Formation (ca.\ 3900--2900~BCE)}\label{sec:naqada}

The Naqada sequence (I: 3900--3650~BCE, II:
3650--3300~BCE, III: 3300--2900~BCE) shows monotonic
increase in all channels.  Nubt (``Gold Town'') becomes
a major center.  Trade networks extend to Mesopotamia.
Art styles show increasing standardization.  By
Naqada~III, kingship iconography and proto-writing
appear.

\textbf{Kernel state at Naqada~III.}  All channels
reach near-maximum: $c_i \in [0.7, 0.95]$.
$F \approx 0.85$, $\omega \approx 0.15$,
$\IC \approx 0.83$, $\Delta \approx 0.02$.
The system enters Watch regime approaching Stable.

\begin{remark}[The ``False Dawn'' of Egyptian
Agriculture]\label{rem:falsedawn}
The archaeological record shows a remarkable
pattern: a crop economy appeared in the Nile Valley
around 7000~BP but was \emph{abandoned}, with
permanent farming delayed to ca.\ 6500~BP.  As
Hillman and Davies demonstrate, crop domestication
itself is rapid (20--200 years), but adoption is not.
In kernel terms, the false dawn has high $F$
(agricultural knowledge exists) but low $\IC$
because the permanent settlement channel $c_3$ is
near-zero (mobile pastoralism persists).  This is
geometric slaughter by the settlement channel.
\end{remark}

\begin{theorem}[The False Dawn as Geometric
Slaughter --- T-NC-4]\label{thm:falsedawn}
The 3rd-millennium false dawn of Egyptian agriculture
(ca.\ 7000~BP) exhibits geometric slaughter:
$F \approx 0.55$ but $\IC \approx 0.12$ because
permanent settlement ($c_3 \to \eps$) destroys
composite integrity while agricultural knowledge
($c_2 \approx 0.7$), material culture ($c_6 \approx
0.6$), and ecological resources ($c_1 \approx 0.8$)
maintain mean fidelity.  $\Delta \approx 0.43$
diagnoses a phase boundary.  Agriculture does not
become structurally real until all channels co-return
($\tR \approx 500$~yr).
\end{theorem}


\subsubsection{Phase~V: Cattle Cult Continuity---Nabta
to Hathor}\label{sec:hathor}

The most striking channel-level continuity spans the
entire 5,000-year arc.  At Nabta Playa (ca.\ 5500~BCE),
cattle were sacrificed in elaborate subterranean
chambers with clay-lined floors and stone roofs.  In
Dynastic Egypt, Hathor---the cow-goddess---is one of
the oldest and most persistent deities, associated
with fertility, the sky, the Milky Way, and the
afterlife.

The structural parallels are precise: bovine ritual
sacrifice $\to$ Hathor temple offerings; stone-chamber
burial $\to$ mastaba and pyramid mortuary architecture;
stellar alignments at Nabta $\to$ stellar religion and
pyramid shaft alignments.  These are not coincidental
similarities---they are the \emph{same channels
persisting through collapse and return}.

\begin{theorem}[Cattle Cult Channel Fidelity ---
T-NC-3]\label{thm:hathor}
The mortuary-ritual channel ($c_4$) maintains
$F > 0.9$ across the Saharan--Nilotic transition
(ca.\ 5500--2500~BCE), a span of $\sim\!3000$~years.
The channel survives the desiccation collapse
(Phase~II) with minimal degradation because ritual
practice is carried by people, not by landscapes.
The Nabta Playa cattle cult $\to$ Hathor cult
trajectory constitutes a single-channel return with
$\tR \approx 0$ (no interruption in ritual practice
despite ecological collapse).
\end{theorem}


% ════════════════════════════════════════════════════════════
\subsection{Egyptian Self-Archaeology as Recursive
Return}\label{sec:selfarch}
% ════════════════════════════════════════════════════════════

A remarkable structural property of Egyptian
civilization---unique in the ancient world---is its
explicit awareness of its own antiquity and its active
efforts to return to earlier states.  This constitutes
a \emph{recursive collapse--return}: the civilization
measuring and re-entering its own collapse history.

\subsubsection{The Evidence}

\textbf{The Dream Stele (ca.\ 1401~BCE).}
Thutmose~IV, as prince, sleeps between the paws of
the Great Sphinx, already buried to its shoulders in
sand.  The Sphinx speaks to him in a dream, promising
the throne in exchange for clearing the sand.  The
stele calls the site the ``Splendid Place of the
Beginning of Time'' (\emph{zep tepi}) and contains a
fragmentary reference to Khaf--- (Khafre), already
1,100+ years in the past.  The stele itself was carved
from stone reused from Khafre's mortuary temple---a
physical re-entry into prior material.

\textbf{Khaemwaset (ca.\ 1250~BCE).}  The fourth son
of Ramesses~II, Khaemwaset is called ``the first
Egyptologist.''  He systematically restored the
pyramids of Unas, Djoser, Sahure, Userkaf, and
Niuserre, adding inscriptions documenting the original
builders.  His inscription reads: ``so greatly did he
love antiquity and the noble folk who were
aforetime.''  He was deliberately measuring and
repairing the drift ($\omega$) accumulated over a
millennium.

\textbf{The Saite Period (664--525~BCE).}  Two
thousand years after the pyramid builders, the 26th
Dynasty deliberately revived Old Kingdom artistic
styles, administrative titles, and religious texts.
This is not nostalgia---it is systematic archaism:
the deliberate return to a prior state vector.

\textbf{The Inventory Stela (ca.\ 670~BCE).}  This
Late Period stela claims that the Temple of Isis
predated the Great Pyramids and that Khufu
``discovered'' an existing temple beside the Sphinx.
Scholarly consensus treats this as a Saite-period
fabrication~\cite{hawass1996sphinx}.  But
structurally, even the \emph{forgery} reveals the
collapse--return pattern: the need to establish
continuity with a deeper past, to demonstrate that
what exists now has always existed.

\textbf{Zep Tepi (``The First Time'').}  The Egyptian
creation concept \emph{zp tpj} describes a primordial
moment when the world emerged from chaos waters (Nu),
a pyramidal mound (benben) rose from the flood, and
the sun appeared.  This imagery maps directly onto the
annual Nile flood---and onto the Saharan--Nilotic
transition itself: water recedes, land appears, life
returns.  The creation myth is the collapse--return
cycle encoded as cosmogony.

\subsubsection{Kernel Analysis}

\begin{theorem}[Self-Archaeology as Recursive Return
--- T-NC-5]\label{thm:selfarch}
Egyptian self-archaeological practice (Khaemwaset
restorations, Dream Stele, Saite archaism) constitutes
a recursive collapse--return: the civilization
operating as its own measurement apparatus, computing
$\omega$ (what has drifted), $F$ (what has persisted),
and $\tR$ (whether return is achievable).  The
recursive return has $\tR_{\mathrm{recursive}} \ll
\tR_{\mathrm{original}}$: Khaemwaset's restorations
span decades, not millennia, because the measurement
infrastructure (writing, administration, religious
continuity) is itself a return channel.
\end{theorem}

\begin{remark}
Four distinct cosmogonies (Hermopolis/Ogdoad,
Heliopolis/Atum, Memphis/Ptah, Thebes/Amun)
share the same structural skeleton: chaos
$\to$ emergence $\to$ order.  In GCD terms, all four
model the same collapse--return cycle with different
Tier-2 channel selections.  The Memphis theology---in
which Ptah creates by speech (``the word'')---is
particularly notable: it identifies the generative
mechanism with the \emph{act of naming}, which is
precisely the Tier-2 operation of channel selection.
\end{remark}


% ════════════════════════════════════════════════════════════
\subsection{Generational Transmission and Drift
Budget}\label{sec:generations}
% ════════════════════════════════════════════════════════════

How do practices persist across 3,000 years without
writing?  And why do some practices stop while others
survive collapse?  The kernel provides a framework for
analyzing generational transmission as a drift budget.

\subsubsection{Transmission Timescales}

At a generation time of $\sim\!25$~years, the
Nabta-to-Dynastic span (5500--3100~BCE) covers
$\sim\!96$ generations.  The pre-Dynastic Nile period
(5000--3100~BCE) covers $\sim\!76$ generations.

The Neolithic Revolution data give concrete bounds.
Hillman and Davies show that crop domestication occurs
in 20--200~years (1--8 generations) once selection
pressure is applied, but \emph{adoption} of
agriculture spreads at 0.6--1.3~km/yr---a demographic
timescale, not a knowledge timescale.

\subsubsection{Why Things Stop: Channel Death Mechanisms}

The kernel identifies four distinct mechanisms by
which channels lose fidelity:

\begin{enumerate}[leftmargin=*]
\item \textbf{Environmental forcing} ($c_1 \to \eps$):
The Saharan desiccation kills ecological
channels regardless of social adaptation.

\item \textbf{Substrate loss}: Practices tied to
specific landscapes (seasonal lake calendrics, desert
megaliths) lose their substrate when the landscape
changes.  The practice ``remembers'' but cannot
re-enact.

\item \textbf{Institutional decay}: Without writing
or rigid institutional structures, high-precision
knowledge degrades at $\sim\!0.5$\% per generation
through oral transmission error.

\item \textbf{Functional displacement}: New
channels displace old ones---agriculture displaces
pure pastoralism, river calendrics displace desert
stellar calendrics, not because the old knowledge is
lost but because new channels are more adaptive.
\end{enumerate}

\subsubsection{Why Things Persist: High-Fidelity
Channels}

Channels that survive collapse share structural
features:

\begin{enumerate}[leftmargin=*]
\item \textbf{Embodied in ritual}: Cattle sacrifice,
burial orientation, solar observation are
performed actions, not stored information.
Performance-based channels resist oral transmission
error because the body remembers.

\item \textbf{Ecologically decoupled}: Ritual does
not require rain.  A cattle sacrifice in the desert
or beside the Nile activates the same channel.

\item \textbf{Socially embedded}: Mortuary practice
is communal and therefore redundantly encoded across
the population.  No single point of failure.
\end{enumerate}

\begin{theorem}[Generational Drift Budget ---
T-NC-6]\label{thm:drift}
Generational habit transmission operates with per-
generation drift $\omega_{\mathrm{gen}} \approx 0.005$
for embodied ritual channels and
$\omega_{\mathrm{gen}} \approx 0.02$ for abstract
knowledge channels.  Over 100 generations:
\beq
\omega_{100} = 1 - (1 - \omega_{\mathrm{gen}})^{100}
\eeq
yields $\omega_{100} \approx 0.39$ for
ritual channels (approaching Collapse) and
$\omega_{100} \approx 0.87$ for abstract channels
(deep Collapse).  The 200-year return time
$\tR \approx 8$ generations represents the
characteristic renewal cycle---the span within which
accumulated drift is corrected through deliberate
re-enactment before fidelity drops below the Watch
threshold.
\end{theorem}


% ════════════════════════════════════════════════════════════
\subsection{The Archaeological Visibility
Problem}\label{sec:visibility}
% ════════════════════════════════════════════════════════════

Between 8000 and 5000~BCE, the Nile world sustained
complex societies (Nabta Playa, Faiyum~A, Merimde,
Badarian), yet the archaeological trace is thin
compared to contemporary Near Eastern sites.  The
kernel explains why.

\textbf{The Saharan erasure.}  The same desiccation
that drove populations to the Nile also buried or
destroyed Saharan sites under sand.  High-fidelity
cultural systems ($F \approx 0.7$) leave low-fidelity
archaeological records because the recording medium
(the landscape) collapses independently of the culture.

\textbf{Perishable media bias.}  Saharan and early
Nilotic societies built in organic materials (wood,
reed, hide, textile) that do not survive 7,000 years
in the archaeological record.  Stone construction
appears late precisely because it was unnecessary in
the Saharan pastoral context---stone was for ritual
(megaliths), not habitation.  The absence of evidence
is structurally predicted by the channel model: the
``material permanence'' channel ($c_6$) was optimized
for function, not for archaeological survival.

\textbf{The world population constraint.}  Global
population in 10,000~BCE was approximately 4~million,
rising to only $\sim$5~million by 5,000~BCE.  The
entire Nile world population during Phase~I--III may
have numbered in the tens of thousands.  Low
population density produces thin archaeological
signal by construction.

\textbf{Agriculture did not help.}  Contrary to
progressive narratives, the adoption of agriculture
\emph{decreased} nutrition and increased disease
burden.  The ``Neolithic Revolution'' was not a
voluntary upgrade but a forced adaptation to
environmental collapse---a return strategy with high
short-term drift cost.  In kernel terms: $F$ of the
subsistence channel initially \emph{drops} during the
agricultural transition before recovering over
several centuries.


% ════════════════════════════════════════════════════════════
\section{Part~II: The Mesopotamian Arc
(5500--2004~BCE)}\label{sec:sumer}
% ════════════════════════════════════════════════════════════

If the Nile world demonstrates collapse--return driven
by environmental forcing (Saharan desiccation), the
Mesopotamian arc demonstrates a structurally distinct
pattern: collapse--return driven by \emph{institutional
overreach}.  The same 7-channel trace vector applies,
but the channel dynamics are inverted.  Where Egypt's
ecological channel died and ritual persisted, Sumer's
ecology was engineered into extraordinary productivity
while social and institutional channels accumulated
fatal drift.

The Sumerians called themselves \emph{sag-gi-ga}
(``black-headed people'') and their land
\emph{ki-en-gi} (``place of the civilized lords'').
Their language---agglutinative, ergative, and without
known relatives---is itself a channel anomaly: a
linguistic isolate that persisted as a scholarly and
liturgical language for two millennia after its last
native speakers died, echoing the Egyptian pattern of
civilizational self-archaeology but through a
different medium.

\subsection{Sumerian Channel Adaptations}\label{sec:sumer-channels}

The 7-channel framework (Table~\ref{tab:channels})
maps to Mesopotamian observables as follows:

\begin{table}[h]
\centering
\caption{Channel proxies adapted for the
Mesopotamian case.  Channels 1--7 are structurally
identical to the Egyptian analysis; the observable
proxies shift to reflect alluvial ecology and cuneiform
documentation.}
\label{tab:sumer-channels}
\begin{tabular}{@{}clP{4.2cm}@{}}
\toprule
$i$ & Channel & Mesopotamian Proxy \\
\midrule
1 & Ecological stability & Irrigation capacity,
    salinization indices, flood regularity \\
2 & Subsistence mode & Cereal yield records,
    barley/wheat ratio (salinization marker) \\
3 & Social complexity & Administrative tablet counts,
    seal usage, settlement hierarchy \\
4 & Mortuary elaboration & Royal tombs, burial
    differentiation, ritual deposits \\
5 & Astronomical knowledge & Astronomical tablets
    (MUL.APIN tradition), calendar reforms \\
6 & Material technology & Bronze metallurgy, wheel,
    pottery mass production, architectural scale \\
7 & Long-distance exchange & Lapis lazuli (Afghanistan),
    carnelian (Indus), timber (Lebanon), copper (Oman) \\
\bottomrule
\end{tabular}
\end{table}


\subsection{Mesopotamian Phases as Kernel States}\label{sec:sumer-phases}

\subsubsection{Phase~S-I: The Ubaid Substrate
(ca.\ 5500--4000~BCE)}\label{sec:ubaid}

The Ubaid period established the template for all
subsequent Mesopotamian civilization.  Originating in
the southern marshlands, Ubaid culture spread across
the entire Fertile Crescent---from the Persian Gulf to
the Mediterranean coast---without evidence of military
conquest.  This is a cultural diffusion event spanning
$\sim$1,500 years.

The Ubaid signature: distinctive painted pottery,
tripartite house plans centered on a T-shaped central
hall, and---crucially---the first monumental temples,
built on platforms that would evolve into
ziggurats~\cite{oates2004ubaid}.  Eridu, later
remembered by Sumerians as the first city, has an
unbroken sequence of 19 temple rebuildings from
Ubaid~I through the Early Dynastic period---each new
temple built directly atop the ruins of its
predecessor.

\textbf{Kernel state.}  Ecological stability is high
in the marshland environment ($c_1 \approx 0.8$).
Subsistence is mixed agriculture-fishing ($c_2
\approx 0.7$).  Social complexity shows temple-centered
organization ($c_3 \approx 0.65$).  Mortuary practice
is moderately elaborate ($c_4 \approx 0.5$).
Astronomical knowledge is implicit in agricultural
calendrics ($c_5 \approx 0.4$).  Material technology
includes fine ceramics but no metals ($c_6 \approx
0.5$).  Exchange networks are extensive (Ubaid pottery
found 800+~km from origin) ($c_7 \approx 0.7$).

$F \approx 0.61$, $\omega \approx 0.39$,
$\IC \approx 0.57$, $\Delta \approx 0.04$.
Watch regime with moderate channel homogeneity.  The
system is structurally comparable to Nabta Playa
Phase~I but with stronger exchange and weaker
astronomical channels.

\begin{remark}
The Eridu temple sequence---19 consecutive
rebuildings on the same spot---is a physical
instantiation of the collapse--return cycle.  Each
destruction event ($\omega \to$ high) is followed by
rebuilding on the exact footprint ($\tR < \IR$),
with the rubble of the prior temple becoming the
foundation platform of the next.  The ziggurat is
literally \emph{accumulated returns}: collapsed
temples raising the next temple closer to the sky.
\end{remark}


\subsubsection{Phase~S-II: The Uruk Explosion
(ca.\ 4000--3100~BCE)}\label{sec:uruk}

The Uruk period represents the most rapid channel
acceleration in the pre-classical world.  Within
900~years, southern Mesopotamia produced: the world's
first city (Uruk, population 40,000--80,000 by
3100~BCE), the first writing system (proto-cuneiform,
ca.\ 3400--3200~BCE), the first monumental architecture
at urban scale (the Eanna precinct and White Temple),
the first mass-produced pottery (the beveled-rim bowl,
interpreted as a standardized ration vessel), and the
first long-distance colonial network (Uruk colonies
from the Nile to the Indus).

Writing emerged not from literature or religion but
from \emph{accounting}---clay tokens sealed in
bullae, then impressed on tablets to track grain
stores and labor obligations.  The earliest tablets
are lists: lists of commodities, lists of professions
(the ``Standard Professions List'' enumerates over
100 occupational categories), lists of cities.  The
Sumerians did not invent writing to tell stories;
they invented it to \emph{audit}.

\textbf{Kernel state.}  All channels spike:
$c_i \in [0.7, 0.95]$.  Ecological stability remains
high via irrigation engineering ($c_1 \approx 0.85$).
Subsistence is maximally productive ($c_2 \approx
0.9$).  Social complexity reaches unprecedented levels
($c_3 \approx 0.95$).  Mortuary elaboration is
moderate relative to social complexity ($c_4 \approx
0.6$).  Astronomical knowledge advances ($c_5 \approx
0.7$).  Material technology includes cylinder seals,
copper alloys, the potter's wheel ($c_6 \approx
0.85$).  Exchange networks span from Egypt to the
Indus ($c_7 \approx 0.9$).

$F \approx 0.82$, $\omega \approx 0.18$,
$\IC \approx 0.78$, $\Delta \approx 0.04$.
Watch regime approaching Stable.

\begin{theorem}[Writing as Ledger Infrastructure ---
T-SM-1]\label{thm:writing}
The invention of cuneiform writing (ca.\ 3400--3200~BCE)
is structurally an integrity ledger: a technology for
tracking channel states (grain stores, labor, trade)
that makes drift $\omega$ \emph{measurable} and
therefore correctable.  Pre-literate societies
accumulate drift silently; literate societies detect
drift via the written record and initiate return before
$\omega$ crosses the Collapse threshold.  Writing
reduces $\tR$ by converting implicit channel states to
explicit, auditable records.
\end{theorem}


\subsubsection{Phase~S-III: Early Dynastic Fragmentation
(ca.\ 2900--2334~BCE)}\label{sec:earlydynastic}

The Early Dynastic period fractures the Uruk unity into
competing city-states: Ur, Uruk, Lagash, Umma, Kish,
Nippur, Eridu.  The \emph{Sumerian King List}---composed
centuries later---retrospectively imposes sequential
dynasties, but the archaeological reality is
simultaneous polities in chronic conflict over water
rights and agricultural land.

This is a structural fragmentation: high overall $F$
but with the social complexity channel ($c_3$)
splitting into \emph{intra-city} coherence (high) and
\emph{inter-city} coherence (low).  The Royal Tombs of
Ur (ca.\ 2600--2400~BCE) demonstrate the extreme:
Queen Puabi's tomb contained 52 ritually sacrificed
attendants, gold and lapis lazuli artifacts of
extraordinary craftsmanship, and the earliest known
stringed instruments.

\textbf{Kernel state.}  Ecological stability begins
its long decline---salinization from irrigation is
detectable in the barley/wheat ratio shift
($c_1: 0.85 \to 0.7$).  Subsistence remains high
($c_2 \approx 0.8$).  Social complexity is paradoxical:
individually high per city but systemically fragmented
($c_3 \approx 0.7$ net).  Mortuary elaboration peaks
($c_4 \approx 0.9$---Ur royal tombs).  Astronomical
knowledge advances ($c_5 \approx 0.7$).  Material
technology peaks in bronze and goldwork ($c_6 \approx
0.9$).  Exchange networks remain extensive ($c_7
\approx 0.8$).

$F \approx 0.79$, $\omega \approx 0.21$,
$\IC \approx 0.77$, $\Delta \approx 0.02$.
Watch regime---high fidelity, low gap, but the
ecological decline ($c_1$ falling) is accumulating
drift that will not manifest for centuries.

\begin{theorem}[Salinization as Slow-Channel
Collapse --- T-SM-2]\label{thm:salinization}
Irrigation-induced salinization constitutes a
slow-channel collapse: $c_1$ degrades at
$\sim$0.003 per generation over $\sim$600 years
(Ubaid through Ur~III), invisible in mean fidelity $F$
until the cumulative drift crosses a phase boundary.
The barley/wheat ratio serves as a proxy:
wheat~(salt-sensitive) disappears from southern
Mesopotamian agriculture by ca.\ 2100~BCE, marking
the moment $c_1$ drops below the Watch threshold.
Unlike the Saharan desiccation (sudden environmental
forcing), salinization is \emph{endogenous}
collapse---the system's own success destroys its
ecological substrate.
\end{theorem}


\subsubsection{Phase~S-IV: Akkadian Unification and
Collapse (ca.\ 2334--2154~BCE)}\label{sec:akkad}

Sargon of Akkad (r.\ ca.\ 2334--2279~BCE) created the
first multi-ethnic empire in recorded history,
unifying the Sumerian south with the Akkadian
(Semitic-speaking) north under centralized
administration.  His grandson Naram-Sin declared
himself a living god---the first ruler to do so.

The Akkadian Empire collapsed within 180~years.
The 4.2~kiloyear climatic event (ca.\ 2200~BCE)---a
severe aridification episode---is widely cited as a
contributing factor, but the kernel analysis reveals a
structural vulnerability: Akkadian centralization
raised $F$ by homogenizing administrative channels
while simultaneously increasing $C$ (curvature---
coupling to uncontrolled degrees of freedom) through
dependence on a single institutional hierarchy.

\textbf{Kernel state at peak (ca.\ 2270~BCE).}
$F \approx 0.80$, $\IC \approx 0.75$, $\Delta
\approx 0.05$.  But $C \approx 0.35$---significantly
higher than any prior phase.  The system is in Watch
regime with elevated curvature, meaning that a
perturbation in \emph{any} channel propagates through
the centralized structure.

\textbf{Kernel state at collapse (ca.\ 2154~BCE).}
The Gutian invasion coincides with the 4.2~ky event.
Ecological channel drops ($c_1: 0.7 \to 0.3$),
administrative channel collapses ($c_3: 0.85 \to
0.2$).  $F \approx 0.48$, $\IC \approx 0.15$,
$\Delta \approx 0.33$.  Collapse regime---and the
$\Delta$ spike is comparable in magnitude to Egypt's
Saharan desiccation ($\Delta \approx 0.32$).

\begin{theorem}[Centralization as Curvature
Amplifier --- T-SM-3]\label{thm:centralization}
Imperial centralization (Akkadian model) increases
mean fidelity $F$ by enforcing channel uniformity
while simultaneously raising curvature $C$---coupling
to uncontrolled degrees of freedom.  When any
channel is perturbed, high $C$ propagates the
perturbation to all other channels through the
centralized administration (single point of
failure).  The Akkadian collapse demonstrates:
$\Delta F / \Delta t$ is $\sim$3$\times$ faster for
the Akkadian collapse ($\sim$50~years) than for the
Saharan desiccation ($\sim$1000~years), because
centralization converts slow environmental drift
into rapid institutional cascade.
\end{theorem}


\subsubsection{Phase~S-V: The Ur~III Renaissance
(ca.\ 2112--2004~BCE)}\label{sec:uriii}

The Third Dynasty of Ur represents Sumer's final and
most bureaucratically complete political expression.
Ur-Nammu (r.\ ca.\ 2112--2095~BCE) produced the oldest
surviving law code---predating Hammurabi by three
centuries---and rebuilt the Great Ziggurat of Ur with
its three-stage platform reaching $\sim$30~meters.

Under his son Shulgi (r.\ ca.\ 2094--2047~BCE), the
Ur~III state created the most extensively documented
bureaucracy in the ancient world: tens of thousands
of administrative tablets record every transaction,
from individual worker rations to international trade
agreements.  Shulgi reformed weights and measures,
standardized the calendar, established scribal
academies (\emph{edubba}, ``tablet houses''), and
composed (or commissioned) hymns declaring himself
divine.

\textbf{Kernel state.}  $F \approx 0.78$ but with an
ominous structure: administrative channels are
near-maximal ($c_3 \approx 0.95$) while the
ecological channel continues its decline ($c_1
\approx 0.55$).  $\IC \approx 0.68$, $\Delta
\approx 0.10$.  The heterogeneity gap is widening:
institutional channels diverge from ecological
reality.

\begin{theorem}[Bureaucratic Overshoot ---
T-SM-4]\label{thm:bureaucracy}
The Ur~III state demonstrates bureaucratic overshoot:
administrative channel fidelity ($c_3 \approx 0.95$)
exceeds ecological reality ($c_1 \approx 0.55$) by
$\sim$0.40.  The administrative apparatus measures
and records channel states with unprecedented
precision but cannot correct the underlying ecological
drift.  This produces a characteristic signature:
$\Delta$ \emph{widens} during the ``renaissance''
because institutional channels accelerate faster
than ecological channels can recover.  The system
appears to be in Watch-approaching-Stable by $F$
alone, but $\Delta$ reveals accumulating structural
strain.
\end{theorem}


\subsubsection{Phase~S-VI: Terminal Collapse
(ca.\ 2004~BCE)}\label{sec:uriii-collapse}

The Ur~III state collapsed in 2004~BCE when Elamite
and Amorite forces sacked Ur.  The \emph{Lament for
Ur}---one of the oldest literary expressions of
civilizational grief---records the destruction:
temples desecrated, canals breached, the population
scattered.  But the kernel reveals that the military
event was an \emph{expression} of structural collapse,
not its cause.

By 2004~BCE, southern Mesopotamian agriculture had
been severely compromised by salinization.  The
population center of gravity had shifted northward
(to Babylon, Isin, Larsa).  The administrative
apparatus that Shulgi built---arguably the most
sophisticated in the ancient world---could not
compensate for ecological channel death.

\textbf{Kernel state.}  $c_1 \approx 0.25$
(environmental collapse in the south), $c_2
\approx 0.4$ (declining yields), $c_3 \to \eps$
(state apparatus destroyed), $c_4 \approx 0.3$
(temples desecrated), $c_5 \approx 0.6$
(astronomical knowledge maintained in scribal
tradition), $c_6 \approx 0.7$ (technology survives),
$c_7 \approx 0.3$ (trade networks disrupted).

$F \approx 0.38$, $\IC \approx 0.06$,
$\Delta \approx 0.32$.  Deep Collapse, with $\Delta$
matching the Egyptian desiccation maximum.

\begin{theorem}[Sumerian as Liturgical Return ---
T-SM-5]\label{thm:sumerian-lang}
Sumerian ceased to be a spoken language by ca.\
2000~BCE but persisted as a scholarly, liturgical,
and literary language for nearly two millennia
thereafter---through the Old Babylonian, Kassite,
Assyrian, and Neo-Babylonian periods, until the
last cuneiform tablets of the Seleucid era (ca.\
1st century~BCE).  This constitutes the longest
documented single-channel return in the ancient
world: linguistic fidelity $F_{\mathrm{ling}} > 0.7$
maintained for $\sim$2,000~years after native-speaker
death, via institutional preservation (scribal
schools) that decoupled the language channel from its
demographic substrate.  The structural parallel to
the Latin persistence in medieval Europe and the
Egyptian cattle-cult survival through desiccation is
exact: channels survive collapse when they are
institutionally or ritually embedded rather than
demographically dependent.
\end{theorem}


\subsection{The Sumerian King List as Recursive
Return}\label{sec:kinglist}

The \emph{Sumerian King List} (composed ca.\
2100~BCE, with later redactions) records kingship
``descending from heaven'' at Eridu before the Flood,
then passing through a sequence of cities after the
Flood.  Pre-Flood kings reign for tens of thousands of
years; post-Flood reigns are centuries, then decades.

This is not na\"{\i}ve mythology.  It is a
\emph{compression algorithm for civilizational
memory}: deep time is represented by extended
reign-lengths, and the Flood serves as the
structural Break---the $\omega \to 1$ event
separating mythic from historical time.  The King
List asserts that kingship \emph{returned} after the
Flood to the same cities, in the same institutional
form.  It is Sumer's equivalent of the Egyptian
\emph{zep tepi} (``the first time'')---a cosmogonic
narrative encoding the axiom's structure.

The Sumerians, like the Egyptians, practiced
self-archaeology.  Nabonidus (r.\ 556--539~BCE),
the last Neo-Babylonian king, excavated foundations
of temples at Ur and Sippar, seeking original
foundation deposits laid 1,500+ years earlier.  His
daughter Ennigaldi-Nanna curated a collection of
artifacts spanning centuries---arguably the world's
first museum.  This is the same recursive return
observed in Khaemwaset: the civilization measuring
its own $\omega$ and attempting $\tR$.

\begin{theorem}[Flood Narrative as Phase-Boundary
Encoding --- T-SM-6]\label{thm:flood}
The Sumerian-Babylonian Flood narrative (Eridu Genesis,
Ziusudra, Atrahasis, Gilgamesh~XI) encodes a
civilizational phase boundary as cosmogony.  The
pre-Flood state (high $F$, low $\omega$, mythic
reign-lengths compressing deep time) transitions
through a total-collapse event ($\omega \to 1$,
all channels $\to \eps$) to a post-Flood return
where kingship ``descends again from heaven'' to
the same cities.  The narrative structure---edenic
state $\to$ catastrophic collapse $\to$ survivor's
return $\to$ restoration---is isomorphic to the
Egyptian \emph{zep tepi} model
(chaos waters $\to$ mound emergence $\to$ solar
return) and both are structural encodings of
Axiom-0: only what returns from collapse is real.
\end{theorem}


% ════════════════════════════════════════════════════════════
\section{Part~III: The Highland Arc
(1000~BCE--1270~CE)}\label{sec:ethiopia}
% ════════════════════════════════════════════════════════════

Ethiopia is Africa's greatest surprise to the
uninitiated---a highland plateau civilization that
developed in geographic isolation, perched above
surrounding desert and lowland swamp, producing a
cultural trajectory with no direct parallel in either
the Nile or Mesopotamian world.  The oldest hominid
fossils (``Ardi,'' \emph{Ardipithecus ramidus},
4.2~Ma; ``Lucy,'' \emph{Australopithecus afarensis},
3.2~Ma) and the oldest stone tools (3.0--3.4~Ma) all
come from Ethiopia~\cite{phillipson2012aksum}.  But
the GCD-relevant arc begins with the first complex
polity at Yeha around 980~BCE and runs through the
Aksumite empire, its Christianization, the Islamic
isolation, and the Zagwe--Solomonic return.

The connection to Part~I is not merely thematic:
Ethiopia is widely identified by scholars with the
``Land of Punt'' described by Egyptian chroniclers as
a trading partner from over 5,000~years ago.  Myrrh,
ivory, animal skins, and gold were imported from Punt
along Nile and Red Sea trade routes.  The highland arc
thus shares a \emph{physical trade channel} with the
Nile world from its earliest documentary
appearance---two civilizations linked by
long-distance exchange while developing independent
collapse--return dynamics on different ecological
substrates.

Where Egypt's ecological channel collapses (Saharan
desiccation) and Sumer's ecological channel slowly
dies (salinization), Ethiopia's highland ecology
provides a \emph{permanent floor}: reliable rainfall,
fertile volcanic soils, and indigenous crops (teff,
ensete, nug) that require no irrigation
infrastructure.  This ecological buffer makes
Ethiopia's collapse dynamics structurally distinct
from both prior cases.

\subsection{Ethiopian Channel
Adaptations}\label{sec:ethiopia-channels}

The 7-channel model requires domain-specific
adaptation for the highland arc:

\begin{table}[h]
\centering
\caption{Ethiopian channel bindings.  The ecological
channel is structurally buffered by highland
altitude.}
\label{tab:ethiopia-channels}
\begin{tabular}{@{}lp{4.2cm}@{}}
\toprule
Channel & Ethiopian Binding \\
\midrule
Ecological stability & Highland rainfall, volcanic
soil, teff/ensete agriculture \\
Subsistence mode & Mixed agro-pastoral (cattle,
plow agriculture, terracing) \\
Social complexity & D'mt $\to$ Aksumite hierarchy
$\to$ Zagwe/Solomonic monarchy \\
Mortuary-ritual & Stelae fields, rock-hewn churches,
Orthodox liturgy \\
Symbolic knowledge & Ge'ez script, Aksumite
coinage, Ethiopian calendar \\
Material technology & Obelisk engineering,
rock-cutting, metalwork, coinage \\
Long-distance exchange & Red Sea trade: India,
Arabia, Rome/Byzantium \\
\bottomrule
\end{tabular}
\end{table}

The critical structural difference: the ecological
and subsistence channels remain at $c_i > 0.7$
across \emph{all} Ethiopian phases.  This creates an
asymmetry with Egypt and Sumer where the ecological
channel is the primary collapse driver.

\subsection{Ethiopian Phases as Kernel
States}\label{sec:ethiopia-phases}

% ── Phase E-I ──
\subsubsection{Phase~E-I: The D'mt Kingdom
(980--400~BCE)}\label{sec:dmt}

The kingdom of D'mt, centered at Yeha in the
northern highlands, represents Ethiopia's first
complex polity.  A Sabaean-style temple at Yeha
(ca.\ 800~BCE) is the oldest standing structure in
sub-Saharan Africa~\cite{phillipson2012aksum}.  The
D'mt polity fuses indigenous highland traditions
with South Arabian cultural elements arriving across
the Red Sea---architectural forms, religious
practice, and possibly early Semitic script.

This is not colonization.  The material culture at
Yeha is predominantly local; South Arabian elements
are selectively adopted into an existing
organizational substrate.  The fusion raises $F$
through channel diversification: trade connections
(new), symbolic systems (enhanced), institutional
forms (hybridized), while the ecological and
subsistence channels continue their highland
baseline.

Kernel state: $F \approx 0.58$, $\IC \approx 0.53$,
$\Delta \approx 0.05$, $C \approx 0.18$.  The
moderate $F$ reflects a nascent polity; the low
$\Delta$ indicates that all channels are developing
roughly in parallel.  Watch regime.

% ── Phase E-II ──
\subsubsection{Phase~E-II: Aksumite Rise
(400~BCE--330~CE)}\label{sec:aksum-rise}

D'mt gives way to the Aksumite Kingdom, which by
the 1st~century~CE sits at a commercial crossroads
linking Egypt, the goldfields of Sudan, the Red Sea,
and the Indian Ocean.  The Persian writer Mani
(216--274~CE) listed Aksum alongside Rome, Persia,
and China as the four great powers of his
era~\cite{munrohay1991aksum}.  Aksum mints its own
coinage in gold, silver, and bronze---the only
sub-Saharan African polity to do so in antiquity.
Monumental stelae up to 33~meters tall mark the
royal burial fields, the tallest monolithic
structures in the ancient world.

The Red Sea trade network drives rapid fidelity
increase: the long-distance exchange channel rises
from $\sim$0.3 (D'mt) to $\sim$0.8 (peak Aksum),
and material technology follows as trade brings
new techniques and materials.  But this very
coupling creates curvature: Aksum's prosperity
depends on external networks it does not control.

Kernel state: $F \approx 0.78$, $\IC \approx 0.72$,
$\Delta \approx 0.06$, $C \approx 0.32$.  High
fidelity, moderate heterogeneity gap, but elevated
curvature from trade dependency.  Watch regime at
upper boundary.

% ── Phase E-III ──
\subsubsection{Phase~E-III: Christian Peak
(330--640~CE)}\label{sec:christian}

Christianity arrives with Frumentius, who converts
King Ezana around 330~CE.  Under Ezana and his
successor Kaleb, the empire reaches its territorial
and commercial peak, extending from modern Somalia
across the Red Sea to the southern Arabian
Peninsula~\cite{munrohay1991aksum}.  Near the end of
the 5th~century, the Nine Saints---monks from Syria,
Greece, and Constantinople---arrive and
systematically establish monasteries throughout the
highlands.

\begin{theorem}[Christianity as Channel Fusion ---
T-ET-1]\label{thm:fusion}
The Christianization of Aksum (330--500~CE) fuses
the mortuary-ritual and symbolic-knowledge channels
into a single high-fidelity compound channel.  The
Ge'ez script, previously used for royal inscriptions
and commercial records, becomes the medium for
liturgical literature, biblical translation (the
Ethiopian Bible includes books absent from other
canons), and monastic scholarship.  Religious
practice and symbolic literacy merge into one
institution (the Ethiopian Orthodox Church).  Unlike
Sumer, where institutional decline fragments
previously coupled channels, this fusion
\emph{raises} $\IC$ by reducing inter-channel
heterogeneity: two separate channels at $c_i
\approx 0.7$ become one compound channel at
$c_{\mathrm{fused}} \approx 0.85$.
\end{theorem}

Kernel state: $F \approx 0.83$, $\IC \approx 0.80$,
$\Delta \approx 0.03$, $C \approx 0.30$.  The
lowest $\Delta$ in the Ethiopian arc: channel fusion
brings the geometric mean close to the arithmetic
mean.  This is the structural peak---high fidelity,
high integrity, concentrated coherence.  But the
curvature remains elevated from trade dependency.

% ── Phase E-IV ──
\subsubsection{Phase~E-IV: Islamic Isolation
(640--1137~CE)}\label{sec:isolation}

The rise of Islam in Arabia (7th~century) disrupts the Red Sea trade
routes that sustained Aksumite prosperity.  The last
Aksumite coins are minted in the 7th~century.
Historians attribute the decline to persistent
drought, overgrazing, deforestation, plague, and
trade-route shifts---but the kernel analysis reveals
a precise mechanism: the high-curvature external
channels (trade, imported technology) collapse
simultaneously while the low-curvature internal
channels (highland agriculture, Christian ritual,
Ge'ez literacy) survive.

The historical record goes nearly silent for
$\sim$500~years.  This is Ethiopia's ``dark age''---
not because civilization ceases, but because the
channels that produce \emph{visible} records (coinage,
monumental construction, foreign diplomatic
correspondence) are precisely the ones that die.

\begin{theorem}[Geographic Buffering Increases
$\Delta$ at Collapse --- T-ET-2]\label{thm:buffer}
The highland plateau maintains the ecological channel
at $c_{\mathrm{eco}} > 0.7$ across all Ethiopian
phases, including the Islamic isolation.  This
prevents the total channel collapse observed in
Saharan desiccation (Egypt, T-NC-1) or terminal
salinization (Sumer, T-SM-2).  However, when
external channels (trade, material technology)
collapse to near $\eps$ while internal channels
remain high, the trace vector becomes \emph{more}
heterogeneous than in cases where all channels fall
together.  The heterogeneity gap $\Delta$ reaches
$\approx 0.35$---\emph{higher} than Egypt
($\sim$0.32) or Sumer ($\sim$0.33) at their
respective collapse boundaries---precisely because
the ecological floor splits the trace vector into
surviving and dead components more extremely.
Geographic buffering does not prevent collapse; it
changes its geometry.
\end{theorem}

Kernel state: $F \approx 0.47$, $\IC \approx 0.12$,
$\Delta \approx 0.35$, $C \approx 0.15$.
$\omega = 0.53$: deep Collapse regime.  Curvature
paradoxically \emph{drops} because the surviving
channels are all internal, low-coupling---the
high-curvature trade channels are dead.  IC crashes
despite F remaining higher than in Egyptian or
Sumerian collapse because the variance between
surviving and dead channels is extreme: geometric
slaughter by channel splitting.

% ── Phase E-V ──
\subsubsection{Phase~E-V: Zagwe Renaissance
(1137--1270~CE)}\label{sec:zagwe}

The Zagwe dynasty emerges around Lalibela around
1137~CE, ending $\sim$500~years of political
fragmentation.  Their signature achievement---the
rock-hewn churches of Lalibela, carved \emph{downward}
from living bedrock rather than built upward from
quarried blocks---constitutes the most structurally
permanent form of architecture in the ancient
world~\cite{phillipson2012aksum}.

\begin{theorem}[Lalibela as Lithic Return ---
T-ET-3]\label{thm:lalibela}
The eleven monolithic churches at Lalibela
(ca.\ 1181--1221~CE) represent a form of
self-archaeology more radical than either Egyptian
pyramid restoration (Khaemwaset) or Mesopotamian
foundation excavation (Nabonidus).  Where the
Egyptians \emph{restored} existing structures and
the Sumerians \emph{excavated} buried ones, the
Zagwe architects carved new sacred structures
\emph{into} geological substrate, rendering them
undetachable from bedrock.  This achieves $\tR$
that approaches geological timescales: a carved-out
church cannot be toppled, burned, looted, or eroded
in any human-relevant timeframe.  The return is
literally lithified---fused into the earth's crust.
\end{theorem}

Kernel state: $F \approx 0.67$, $\IC \approx 0.61$,
$\Delta \approx 0.06$, $C \approx 0.20$.  The
return is genuine but partial: ritual and
construction channels are restored to high fidelity,
but long-distance trade remains suppressed.  The
pattern parallels Egypt's Phase~III (Nilotic
consolidation): internal channels drive recovery
while external channels lag.

% ── Phase E-VI ──
\subsubsection{Phase~E-VI: Solomonic
Restoration
(1270--1529~CE)}\label{sec:solomonic}

In 1270, Yekuno Amlak deposes the last Zagwe king
and establishes the Solomonic dynasty, claiming
descent from King Solomon and the Queen of Sheba
through Menelik~I.  The legitimizing text is the
\emph{Kebra Nagast} (``Glory of Kings''), compiled
in its current form around 1300~CE.

\begin{theorem}[The Kebra Nagast as Genealogical
Return --- T-ET-4]\label{thm:kebra}
The \emph{Kebra Nagast} constructs a direct lineage
from Solomon (ca.\ 970~BCE) through Menelik~I to
Yekuno Amlak (1270~CE), asserting $\tR \approx
2{,}240$ years---the longest genealogical return
claim in the documentary record.  This is not
merely mythic legitimation.  The text functions as
a \emph{frozen contract}: it asserts that the
Solomonic line constitutes the only legitimate
channel of royal authority, and this assertion is
never subsequently renegotiated (the dynasty
persists until 1974).  The Kebra Nagast converts
a genealogical claim into a structural invariant:
once frozen, the Solomonic identity becomes the
fixed point around which all subsequent political
dynamics orbit~\cite{budge1922kebra}.
\end{theorem}

Ethiopia in this period again expands militarily to
dominate the Horn of Africa.  The arts and
literature advance.  The empire operates without a
fixed capital---emperors move in mobile camps
across the territory, a practice that distributes
royal presence rather than concentrating it.

Kernel state: $F \approx 0.73$, $\IC \approx 0.68$,
$\Delta \approx 0.05$, $C \approx 0.25$.  Not as
high as the Aksumite peak ($F = 0.83$) because the
Red Sea trade never fully recovers---Muslim control
of maritime routes persists.  But the return is
structurally complete: all seven channels are active,
with no channel near $\eps$.

\subsection{Ge'ez and the Sacerdotal
Parallel}\label{sec:geez}

Ge'ez, the classical Ethiopic language, provides the
most direct structural parallel between the
Ethiopian and Sumerian arcs.  Like Sumerian, Ge'ez
ceases as a spoken vernacular (around the
9th--10th~century~CE) but survives as the liturgical
and scholarly language of the Ethiopian Orthodox
Church---and continues in this role to the present
day.

\begin{theorem}[Ge'ez as Sacerdotal Persistence ---
T-ET-5]\label{thm:geez}
Ge'ez, like classical Sumerian after the fall of
Ur~III (T-SM-5), persists as a non-vernacular
prestige language embedded in ritual and scholarly
institutions for over 1{,}100~years (and counting).
This constitutes single-channel return:
$c_{\mathrm{symbolic}} > 0.7$ while the spoken
language channel transitions to Amharic, Tigrinya,
and other Semitic successors.  The structural
parallel is exact: institutional embedding (temple
in Sumer, church in Ethiopia) preserves a
high-fidelity linguistic channel after the
demographic base has shifted to successor
languages.  Sacerdotal persistence is the general
mechanism; the specific institutions vary.
\end{theorem}

\subsection{Adwa and the Military Return
Channel}\label{sec:adwa}

A structural feature unique to Ethiopia among the
three civilizations is the survival of a military
return channel across the modern threshold.  At the
Battle of Adwa (1~March 1896), Emperor Menelik~II's
forces decisively defeated the Italian colonial army
in what was the first time an African army defeated
a European colonial invasion in open
battle~\cite{marcus2002ethiopia}.  This is not
merely a military event.  In kernel terms, Adwa
constitutes a \emph{return through a channel that
had been declared dead} by the colonial framework:
every other African civilization had its military
channel reduced to $\eps$ by European industrial
warfare.

\begin{theorem}[Adwa as Modern Return ---
T-ET-6]\label{thm:adwa}
The Ethiopian victory at Adwa (1896) constitutes
the only instance in the dataset of a civilization
executing a successful return through the military
channel against an industrial-era opponent.  This is
structurally enabled by the highland buffer
(T-ET-2): because the ecological and subsistence
channels never collapsed, the demographic and
organizational base for mass military mobilization
persisted through the Islamic isolation and
Solomonic periods.  In Egypt and Sumer, the
equivalent military channel had been extinct for
millennia by the colonial era.  Adwa demonstrates
that the highland fidelity floor has cascading
consequences across geological-to-modern timescales.
\end{theorem}


% ════════════════════════════════════════════════════════════
\section{Discussion}\label{sec:discussion}
% ════════════════════════════════════════════════════════════

\subsection{Comparative Kernel Trajectories}

The three arcs, when read through the kernel, reveal
a shared collapse--return architecture with three
distinct ecological substrates and three different
modes of collapse.

\begin{table}[h]
\centering
\caption{Comparative kernel states across Egyptian,
Sumerian, and Ethiopian phases.  $\Delta = F - \IC$
is the heterogeneity gap.}
\label{tab:comparative}
\begin{tabular}{@{}lcccc@{}}
\toprule
Phase & $F$ & $\IC$ & $\Delta$ & Regime \\
\midrule
\multicolumn{5}{l}{\emph{Egypt}} \\
\quad I: Nabta Playa & 0.71 & 0.68 & 0.03 & Watch \\
\quad II: Desiccation & 0.39 & 0.07 & 0.32 & Collapse \\
\quad III: Badarian & 0.66 & 0.63 & 0.03 & Watch \\
\quad IV: Naqada~III & 0.85 & 0.83 & 0.02 & Watch \\
\midrule
\multicolumn{5}{l}{\emph{Sumer}} \\
\quad S-I: Ubaid & 0.61 & 0.57 & 0.04 & Watch \\
\quad S-II: Uruk & 0.82 & 0.78 & 0.04 & Watch \\
\quad S-III: Early Dyn. & 0.79 & 0.77 & 0.02 & Watch \\
\quad S-IV: Akkad peak & 0.80 & 0.75 & 0.05 & Watch \\
\quad S-IV: Akkad fall & 0.48 & 0.15 & 0.33 & Collapse \\
\quad S-V: Ur~III & 0.78 & 0.68 & 0.10 & Watch \\
\quad S-VI: Ur fall & 0.38 & 0.06 & 0.32 & Collapse \\
\midrule
\multicolumn{5}{l}{\emph{Ethiopia}} \\
\quad E-I: D'mt & 0.58 & 0.53 & 0.05 & Watch \\
\quad E-II: Aksumite & 0.78 & 0.72 & 0.06 & Watch \\
\quad E-III: Christian & 0.83 & 0.80 & 0.03 & Watch \\
\quad E-IV: Isolation & 0.47 & 0.12 & 0.35 & Collapse \\
\quad E-V: Zagwe & 0.67 & 0.61 & 0.06 & Watch \\
\quad E-VI: Solomonic & 0.73 & 0.68 & 0.05 & Watch \\
\bottomrule
\end{tabular}
\end{table}

Four structural observations emerge from the
three-way comparison:

\textbf{1.~$\Delta$ at collapse is universal.}  All
three civilizations hit $\Delta \in [0.32, 0.35]$ at
their respective collapse boundaries---despite
completely different ecological substrates, channel
dynamics, and collapse mechanisms.  Egypt collapses
from exogenous environmental forcing; Sumer from
endogenous salinization plus invasion; Ethiopia from
exogenous trade-route disruption.  The heterogeneity
gap is channel-agnostic: it measures structural
divergence regardless of which channel dies.

\textbf{2.~Inverted vulnerability profiles.}  Each
civilization's high-fidelity survival channels are
precisely the ones that the other two lose first.
Egypt's durable channels are ritual and astronomical
(ecologically decoupled).  Sumer's are institutional
and technological (ecologically coupled).  Ethiopia's
are ecological and subsistence (geographically
buffered).  This three-way inversion suggests that
collapse-return dynamics depend not on the specific
channels but on the heterogeneity \emph{structure} of
the trace vector at the moment of regime transition.

\textbf{3.~The highland buffer paradox.}
Ethiopia's ecological floor ($c_{\mathrm{eco}} > 0.7$)
does not prevent Collapse regime entry---it produces
the \emph{highest} $\Delta$ at collapse (0.35 vs.\
0.32--0.33 for Egypt and Sumer).  This is because
geographic buffering splits the trace vector more
extremely: surviving channels remain high while dead
channels hit $\eps$, maximizing the arithmetic-geometric
mean gap.  Buffering changes the geometry of collapse,
not its inevitability.

\textbf{4.~Self-archaeology is universal and
independent.}  All three civilizations---separated by
thousands of kilometers and speaking mutually
unintelligible languages---independently develop
recursive return practices: Khaemwaset restores
pyramids, Nabonidus excavates foundations, and the
Zagwe carve churches into bedrock.  The escalating
permanence (restoration $\to$ excavation $\to$
lithification) suggests a structural gradient: as
$\omega_{\mathrm{accumulated}}$ increases across
civilizational timescales, the return architectures
become increasingly immovable.

\subsection{Generative Collapse: What Returns Is
New}

In all three cases, the return is not a restoration
of the prior state.  Egypt's Nilotic return
(Phase~III) exceeds Nabta Playa (Phase~I) in every
channel.  Sumer's post-Akkadian return (Ur~III)
exceeds the pre-Akkadian Early Dynastic period in
administrative sophistication.  Ethiopia's Zagwe
return (Phase~E-V) produces the rock-hewn churches of
Lalibela---a form of monumental architecture that
surpasses anything built at Aksum in structural
permanence.  The axiom is precise: collapse is
\emph{generative}.  The structure that returns is
richer than what entered the collapse---because the
collapse itself selects for high-fidelity channels
while eliminating low-fidelity ones.

\subsection{Why the Monuments Outlive the
Civilizations}

When Khaemwaset restored the Pyramid of Unas around
1250~BCE, he was looking back 1,300~years.  When the
Dream Stele was carved around 1401~BCE, the Sphinx
was already 1,100+~years old and buried in sand.
When Nabonidus excavated at Ur around 545~BCE, he
sought foundation deposits from $\sim$1500~BCE.
When the Zagwe commissioned the churches at Lalibela
around 1200~CE, they carved \emph{into} the same
geological substrate that predates all human
construction---achieving $\tR$ bounded only by
erosion rates.
These are not isolated antiquarian acts.  They are
civilizations repeatedly measuring their own $\omega$
and computing $\tR$---asking whether their past can
be reached.  The answer was structurally yes, because
high-fidelity channels (ritual, writing, institutional
memory) provided the return path.  The escalating
permanence across the three arcs---Egypt builds
\emph{on} bedrock, Sumer builds \emph{with} mud brick
(and it erodes), Ethiopia carves \emph{into}
bedrock---reflects an increasing sophistication of
the return strategy itself.
\emph{Solum quod redit, reale est}---and all three
civilizations made their past real by returning to it,
through progressively more durable media.

\subsection{Genetics and Continuity}

The 2025 whole-genome analysis of Early Dynastic
Egyptian remains~\cite{marshall2025genome} showing
77.6\% North African Neolithic + 22.4\% eastern
Fertile Crescent ancestry confirms biological
continuity across the Saharan--Nilotic transition:
no population replacement.  In Mesopotamia, genetic
evidence is sparser, but the Ubaid-to-Uruk cultural
continuity at sites like Eridu (19 consecutive temple
rebuildings) and the absence of destruction layers
between periods argue for demographic continuity
with Semitic-speaking Akkadian admixture.

Ethiopian genetic continuity is among the most
robust in the dataset: highland populations show
deep indigenous ancestry with limited external
admixture, consistent with geographic isolation
over millennia~\cite{marcus2002ethiopia}.
The Solomonic dynasty's explicit genealogical
claim (Solomon $\to$ Menelik~I $\to$ Yekuno Amlak)
operates as a \emph{cultural} fidelity channel
independent of the genetic evidence, but the
demographic continuity of the highland population
provides the substrate on which this cultural
claim is sustained.  In all three cases, the
genetic channel has $F > 0.7$ across millennia.


% ════════════════════════════════════════════════════════════
\section{Conclusion}\label{sec:conclusion}
% ════════════════════════════════════════════════════════════

The GCD kernel provides a formal instrument for
analyzing civilizational dynamics that conventional
historiography treats as progressive narrative.
Applied comparatively to the Nile world
(8000--2900~BCE), Mesopotamia (5500--2004~BCE),
and the Ethiopian highlands (1000~BCE--1270~CE),
it reveals a shared collapse--return architecture
operating across three distinct ecological substrates
with three structurally inverted vulnerability
profiles.

Eighteen theorems span the three arcs.  From Egypt:
(T-NC-1)~environmental collapse produces geometric
slaughter; (T-NC-2)~return is measurable through
$\Delta$ closure; (T-NC-3)~cattle-cult fidelity
persists 3,000~years through ecological collapse;
(T-NC-4)~the agricultural false dawn is geometric
slaughter by the settlement channel;
(T-NC-5)~Egyptian self-archaeology is recursive
return; (T-NC-6)~generational drift budgets explain
the 200-year renewal cycle.

From Sumer: (T-SM-1)~cuneiform writing is a ledger
infrastructure that makes drift measurable;
(T-SM-2)~salinization is slow endogenous channel
collapse; (T-SM-3)~centralization amplifies curvature
and converts slow drift to rapid cascade;
(T-SM-4)~bureaucratic overshoot widens $\Delta$
during apparent renaissance; (T-SM-5)~Sumerian
language persistence is the longest documented
single-channel return; (T-SM-6)~the Flood narrative
encodes a phase boundary as cosmogony.

From Ethiopia: (T-ET-1)~Christianization fuses
mortuary-ritual and symbolic channels, raising $\IC$
by reducing heterogeneity; (T-ET-2)~geographic
buffering increases $\Delta$ at collapse rather than
preventing it; (T-ET-3)~the Lalibela rock-hewn
churches are lithified return with geological-scale
$\tR$; (T-ET-4)~the Kebra Nagast freezes a
genealogical return claim across 2,240~years;
(T-ET-5)~Ge'ez sacerdotal persistence parallels
Sumerian as institutional single-channel return;
(T-ET-6)~Adwa constitutes military-channel return
through a threshold declared extinct by the colonial
framework.

The heterogeneity gap $\Delta = F - \IC$ remains
the primary diagnostic across all three
civilizations.  All hit $\Delta \in [0.32, 0.35]$
at collapse---Egypt from exogenous environmental
forcing, Sumer from endogenous institutional
cascade, Ethiopia from exogenous trade disruption.
All independently develop recursive self-archaeology
with escalating permanence.  All encode
collapse--return as cosmogonic or genealogical
narrative (\emph{zep tepi}, the Flood, the
Solomonic lineage).

The Ethiopian arc adds a structural dimension
absent from the Nile--Euphrates comparison: the
persistence of the ecological floor changes the
\emph{geometry} of collapse without preventing its
occurrence.  Geographic buffering splits the trace
vector more extremely (higher $\Delta$) but
preserves the substrate from which return can
operate.  This three-way comparison suggests that
civilizational collapse--return dynamics are
governed by the heterogeneity structure of the
trace vector rather than by the specific ecological,
institutional, or cultural content of any
individual channel.

Where conventional analysis sees emergence,
development, or decline, the kernel sees collapse,
return, and the generative structure that only
becomes visible when a civilization crosses back
through its own phase boundary.

\emph{Finis, sed semper initium recursionis.}

% ════════════════════════════════════════════════════════════
% bibliography
% ════════════════════════════════════════════════════════════
\bibliography{Bibliography}

\end{document}
